\documentclass[a4paper]{article}
\usepackage{graphicx}
\usepackage{amsmath,amssymb}
\usepackage{hyperref}
\usepackage{minted}
\usepackage{multirow}

\title{Adverb}
\date{}
\begin{document}
    \maketitle
    
    
    \section{Modaladverbien (Art und Weise)}
        
        Modaladverbien beschreiben, wie eine Handlung ausgeführt wird.\\
        Beispiel: Er spricht \textbf{schnell}.\\
        Meaning: He speaks \textbf{quickly}.
    
    
    \section{Temporaladverbien (Zeit)}
        
        Temporaladverbien geben an, wann etwas passiert.\\
        Beispiel: Ich komme \textbf{morgen}.\\
        Meaning: I will come \textbf{tomorrow}.
    
    
    \section{Lokaladverbien (Ort)}
        
        Lokaladverbien zeigen den Ort einer Handlung.\\
        Beispiel: Wir bleiben \textbf{hier}.\\
        Meaning: We stay \textbf{here}.
    
    
    \section{Kausaladverbien (Grund)}
        
        Kausaladverbien drücken einen Grund oder eine Folge aus.\\
        Beispiel: Er ist \textbf{deshalb} zu Hause.\\
        Meaning: He is at home \textbf{therefore}.
    
    
    \section{Gradadverbien (Grad / Intensität)}
        
        Gradadverbien verstärken oder schwächen eine Aussage.\\
        Beispiel: Das ist \textbf{sehr} wichtig.\\
        Meaning: That is \textbf{very} important.
    
    
    \section{Konjunktionaladverbien (Satzverknüpfung)}
        
        Konjunktionaladverbien verbinden Sätze und zeigen eine logische Beziehung.\\
        Beispiel: Es regnet. \textbf{Trotzdem} gehe ich spazieren.\\
        Meaning: It is raining. \textbf{Nevertheless} I go for a walk.
    
        \begin{itemize}
            \item deswegen
            \item deshalb
        \end{itemize}

\end{document}